% Volume VI: Adversaries, Platforms, and Automated Judgment
% Part XI. Untrustworthy Actors
% Chapter 59: Models of Adversarial Agency
% Spine: proof-spines/ch059-spine.md

\chapter{Models of Adversarial Agency}
\label{ch:ch059}


Disciplined paranoia begins by replacing generic talk of ``security'' with explicit attacker models of increasing capability. The relevant question is not whether an institution faces adversaries in the abstract, but which forms of agency it can actually resist. The outline therefore orders adversarial power by a nested family
\[
\mathcal A_0\subset\mathcal A_1\subset\cdots\subset\mathcal A_k.
\]
\Definition\ Nested adversary classes order hostile capability by increasing powers.
Each step in the chain adds capacities an actor may possess: to lie, omit, frame, flood, intimidate, collude, impersonate, or exploit procedural rules.

Once these adversary classes are stated, suspicious appearances can be interpreted proportionately. The same surface anomaly may be consistent with error, routine incentive, or intentional manipulation; the institution should say whether its procedures are robust against, say, \(\mathcal A_2\) but not \(\mathcal A_5\), rather than pretending to a universal immunity it does not have. The chapter's preface claim is therefore simple: disciplined suspicion starts with a bounded model of adversarial agency and asks which capabilities must be defeated before judgment can count as secure.
